\subsection*{Walking through the kernel}

The stage index is \co@iv_k % STAGES@. With two slots, tiles 0 and 2 use slot 0, while tiles 1 and 3 use slot 1. These are different logical tiles sharing physical storage at different times; the event pair belongs to the slot, not to one permanent tile.

The producer waits on \co@empty[stage]@ before writing either operand into that slot. Both copies in this example are synchronous, so triggering \co@full[stage]@ after them means both operands are ready. Replacing the copies with asynchronous operations would require completing those operations before sending the same signal.

The consumer first triggers every \co@empty@ event. All events begin unset, so this bootstrap grants the producer permission to fill each slot once. No \co@full@ event is triggered here because no input has been loaded yet.

The consumer waits on \co@full[stage]@, loads the operand fragments, and accumulates the slice. It then performs \co@mma.commit@ before triggering \co@empty[stage]@. That order matters: issuing a matrix operation is not sufficient evidence that it has stopped reading the shared operands.

The completion point protects buffer reuse. If the producer overwrites a slot while MMA still reads it, the consumer can mix operands from two K slices even though the tile shapes are correct. Checking the \co@full@ signal alone cannot detect this mistake; the return of the \co@empty@ credit is equally necessary.
